\chapter{Аномальный индекс и тёплицев коэффициент: внешний продукт или ручная вставка}
\label{ch:v6-spectral-transition-anomaly-to-toeplitz-product-map-gate}

\section{Две карты из одного представительного носителя}

Предыдущая глава показала, что стандартный сфалеронный поток равен
четырём на поколение, тогда как проектный коэффициентный класс имеет
ранг пятнадцать. Настоящий гейт проверяет последнюю естественную
возможность: может ли каспаровский или индексный внешний продукт
автоматически превратить единичный сфалеронный класс в класс пятнадцать.

Относительно слабой группы носитель одного поколения раскладывается как
\begin{equation}
 [H_{15}]\big|_{SU(2)_L}
 =4[\mathbf2]+7[\mathbf1].
 \label{eq:v6-anomaly-product-representation-ring}
\end{equation}
Здесь четыре фундаментальных дублета суть три цветовые копии $Q_L$ и
один $L_L$, а семь синглетов суть $u_R^r,u_R^g,u_R^b$,
$d_R^r,d_R^g,d_R^b,e_R$.

На этом же элементе кольца представлений действуют две разные аддитивные
карты.

\subsection*{Забывающая размерность}

Обычная комплексная размерность удовлетворяет
\begin{equation}
 \dim(\mathbf2)=2,
 \qquad
 \dim(\mathbf1)=1,
\end{equation}
поэтому
\begin{equation}
 \dim H_{15}=4\cdot2+7\cdot1=15.
 \label{eq:v6-anomaly-product-dimension-fifteen}
\end{equation}

\subsection*{Сфалеронное индексное спаривание}

Для единичного $SU(2)$-перехода фундаментальный левый дублет даёт одно
ориентированное пересечение, а слабый синглет --- ноль:
\begin{equation}
 \mathcal I_{\rm sph}(\mathbf2)=1,
 \qquad
 \mathcal I_{\rm sph}(\mathbf1)=0.
\end{equation}
Следовательно,
\begin{equation}
 \mathcal I_{\rm sph}(H_{15})=4\cdot1+7\cdot0=4.
 \label{eq:v6-anomaly-product-index-four}
\end{equation}
Индекс зависит от представления twisting-пучка, а не только от его
ранга; это является стандартным содержанием индексного спаривания и его
эквивариантной формы
\cite{anomaly-product-atiyah-singer1968,
anomaly-product-aps1975}.

\section{Формальный продукт действительно даёт пятнадцать}

Пусть $x_{\rm sph}$ обозначает класс единичного пересечения, а
$[q_0]\in K_0(\mathbb C)$ --- забытый до обычного конечномерного
векторного пространства проектор ранга пятнадцать. Тогда
мультипликативность индекса даёт
\begin{equation}
 \operatorname{Ind}(x_{\rm sph}\boxtimes[q_0])
 =\operatorname{Ind}(x_{\rm sph})\rank(q_0)=1\cdot15=15.
 \label{eq:v6-anomaly-product-formal-fifteen}
\end{equation}
Такой внешний продукт математически допустим в $KK$-языке
\cite{anomaly-product-kasparov1981}. Но он описывает пятнадцать
одинаковых копий одного и того же активного оператора. Он забыл, что
семь состояний являются слабым синглетами, а восемь компонент собраны в
четыре дублета.

Иными словами, формула \eqref{eq:v6-anomaly-product-formal-fifteen}
получает пятнадцать ровно потому, что $[q_0]$ с этим рангом был введён
до спаривания. Она не выводит $q_0$ из сфалеронного гамильтониана.

\section{Эквивариантный продукт возвращает четыре}

Если сохранить физическое $SU(2)_L$-действие, коэффициентом является не
безразличный $\mathbb C^{15}$, а элемент
$4[\mathbf2]+7[\mathbf1]$ из \eqref{eq:v6-anomaly-product-representation-ring}.
Тогда эквивариантное индексное спаривание обязано дать
\begin{equation}
 x_{\rm sph}\widehat\otimes_{SU(2)}
 \left(4[\mathbf2]+7[\mathbf1]\right)
 \longmapsto 4,
\end{equation}
что совпадает с физическим спектральным потоком и аномальными весами
\cite{anomaly-product-klinkhamer-rupp2003}.

Различие нельзя устранить общей нормировкой. Для проектных компонент
\begin{align}
 \dim H_q&=12, & \mathcal I_{\rm sph}(H_q)&=3,\\
 \dim H_\ell&=3, & \mathcal I_{\rm sph}(H_\ell)&=1.
\end{align}
Отношения равны $12/3=4$ и $3/1=3$. Поэтому не существует одного
множителя, переводящего физические компонентные потоки $(3,1)$ в
проектные ранги $(12,3)$.

\begin{theorem}[Формальное и физическое произведения различны]
Неэквивариантное произведение единичного сфалеронного класса с заранее
заданным проектором ранга пятнадцать имеет индекс пятнадцать. Физическое
эквивариантное спаривание того же пути с фактическим носителем одного
поколения $4[\mathbf2]+7[\mathbf1]$ имеет индекс четыре. Следовательно,
число пятнадцать в первом произведении является входным коэффициентом,
а не результатом сфалеронной динамики.
\end{theorem}

\section{Real/KO6-версия}

Для сопряжённых ориентаций формальное декорирование даёт пару
\begin{equation}
 (-15,+15),
\end{equation}
а физическое слабое спаривание ---
\begin{equation}
 (-4,+4).
\end{equation}
Обе пары имеют нулевой обычный суммарный комплексный индекс, но это не
делает их одним Real-классом. Для такого отождествления требовался бы
морфизм коэффициентных алгебр, сохраняющий представления, градуировку и
Real-структуру. Текущий проект такого морфизма не содержит.

\section{Физический смысл отрицательного результата}

Сфалеронный маршрут не опровергает класс пятнадцать как доказанный
тёплицев коэффициент. Он опровергает более сильное утверждение, будто
стандартный электрослабый сфалерон \emph{рождает} все пятнадцать
фермионных направлений одним спектральным потоком.

Сохраняются два разных результата:
\begin{enumerate}
 \item $15$ --- топологический коэффициент и ledger полного пакета
 одного поколения;
 \item $4$ --- физический хиральный поток слабых дублетов в стандартном
 сфалеронном переходе.
\end{enumerate}
Их нельзя смешивать без новой динамики, действующей также на правые
синглеты.

\begin{nogocriterion}
Внешний продукт не считается механизмом рождения класса пятнадцать,
если после забывания gauge-действия он лишь умножает единичный индекс на
заранее известный ранг $q_0$. Физически допустимый продукт обязан
сохранять слабые представления и воспроизводить тот же коэффициент без
замены $4[\mathbf2]+7[\mathbf1]$ на пятнадцать одинаковых копий.
\end{nogocriterion}

\section{Следующая архитектурная развилка}

Следующий гейт
\path{version6_spectral_transition_class15_physical_role_branch_decision_gate}
должен окончательно выбрать роль класса пятнадцать:
\begin{enumerate}
 \item оставить его классификационным ledger пакета поколения, не
 объявляя multiplicity рождения;
 \item либо ввести явно новую переходную динамику полного конечного
 оператора, действующую на дублеты и синглеты и проходящую отдельные
 тесты аномалий, KO6 и устойчивости.
\end{enumerate}

Стоп-критерий: запрещено продолжать стандартный сфалеронный маршрут под
новым названием, если оператор по-прежнему имеет физический индекс
четыре.

\begin{protocol}
\begin{enumerate}
 \item физическое слабое представление: \textbf{$4[\mathbf2]+7[\mathbf1]$};
 \item забывающая размерность: \textbf{$15$};
 \item сфалеронное индексное спаривание: \textbf{$4$};
 \item формальный продукт с забытым $q_0$: \textbf{$15$};
 \item физический эквивариантный продукт: \textbf{$4$};
 \item компонентные ранги: \textbf{$(12,3)$};
 \item компонентные потоки: \textbf{$(3,1)$};
 \item единый нормировочный множитель: \textbf{не существует};
 \item формальная Real-пара: \textbf{$(-15,+15)$};
 \item физическая Real-пара потока: \textbf{$(-4,+4)$};
 \item стандартный сфалерон выводит класс пятнадцать: \textbf{нет};
 \item класс пятнадцать как классификационный ledger: \textbf{сохранён};
 \item физическое замыкание: \textbf{не достигнуто}.
\end{enumerate}
\end{protocol}

Машинный сертификат сохранён в
\path{s2t_v6_spectral_transition_anomaly_to_toeplitz_product_map_gate_results.json}.

\begin{thebibliography}{9}
\bibitem{anomaly-product-atiyah-singer1968}
M.~F.~Atiyah, I.~M.~Singer,
\emph{The Index of Elliptic Operators: I},
Annals of Mathematics 87 (1968), 484--530.

\bibitem{anomaly-product-aps1975}
M.~F.~Atiyah, V.~K.~Patodi, I.~M.~Singer,
\emph{Spectral Asymmetry and Riemannian Geometry I},
Mathematical Proceedings of the Cambridge Philosophical Society 77
(1975), 43--69.

\bibitem{anomaly-product-kasparov1981}
G.~G.~Kasparov,
\emph{The Operator $K$-Functor and Extensions of $C^*$-Algebras},
Mathematics of the USSR-Izvestiya 16 (1981), 513--572.

\bibitem{anomaly-product-klinkhamer-rupp2003}
F.~R.~Klinkhamer, C.~Rupp,
\emph{Sphalerons, Spectral Flow, and Anomalies},
Journal of Mathematical Physics 44 (2003), 3619--3639;
arXiv:hep-th/0304167.
\end{thebibliography}